% ============================================================
% DRAFT NOTE (method)
% Contains: The per-compressor scalar prediction target, causal input policy,
% and learned and non-learned baselines.
% Written now: The settled model family and feature-access constraints.
% Pending: Exact fixed XGBoost hyperparameters and feature inventory.
% Sources: docs/methodology.md; docs/goal.md; docs/_why/3_measuring_quality.md
% ============================================================
\section{Pre-switch Forecaster}

For each compressor \compressor, we fit one scalar regressor
\begin{equation}
\hat d_c(s,r) = f_c(x_{\leq s}, r, s),
\end{equation}
where $x_{\leq s}$ contains only information available before compression is
activated. Supported ratios are pooled within a compressor, with \ratio as an
input. We fit a separate model for each compressor.

The input whitelist contains the known removal ratio, absolute generated
position, and causal pre-switch \texttt{feat\_\_*} statistics. Forbidden
inputs are final reference length, relative position derived from final
length, reference or hybrid quality, labels or label-derived fields, prompt
identity, probes, hybrid-stream summaries, and every other post-switch signal.
Feature timing at \switchpos is validated explicitly to exclude an off-by-one
leak. Figure~\todo{F3 reference} will show this boundary.

The primary learned model is a fixed XGBoost squared-error regressor
\citep{xgboost}. A linear Huber regressor on the same whitelisted inputs is a
learned baseline. Training-only non-learned baselines are the global mean and
median, mean and median by ratio, and mean and median by ratio and 16-token
absolute-position bucket. \todo{Specify the frozen XGBoost hyperparameters
and the exact feature inventory.}
